\section{Affine counting with adaptive auxiliary primes}
\label{sec:affine-counting}

The geometric arguments below place the points to be counted on varieties
whose equations vary with the box.  We first prove a counting theorem
that is uniform for these moving varieties.  Its input is a lower bound
for the number of independent monomials on each subvariety.  Its proof
successively cuts residue classes by polynomial equations.  A deletion
argument controls the primes at which these newly constructed equations
have singular reduction.
The local prime-adic determinant principle is due to Heath-Brown
\cite[Section~3, Lemmas~5--6]{HeathBrown2002Density}; the ordered local Hilbert
and divisibility estimates are developed in
\cite[Lemmas~2.3--2.5]{Salberger2007}. Simultaneous prime
conditions appear in \cite[Theorem~3.2]{Salberger2007} and
\cite[Section~2, Theorem~2.2]{Salberger2023}.
The argument here includes
the additional patch construction and deletion estimate needed when
later varieties depend on earlier prime choices.

Throughout this section, degree means the ordinary degree of projective
closure; for a reducible algebraic set we use the sum of the degrees of
its irreducible components.  Irreducibility is over $\C$.  We use the
degree bounds for intersections and linear projections supplied by
B\'ezout's theorem.  The ambient dimension is fixed.

Fix positive weights $\omega_1,\ldots,\omega_N$.  The weight of the
coordinate monomial $z^\alpha$ is $\omega\cdot\alpha$.  For an affine
algebraic set $A\subseteq\C^N$, define
\[
 H_A(T)=\dim_{\C}\operatorname{span}
 \{z^\alpha|_A:\omega\cdot\alpha\le T\}.
\]
Thus $H_A(T)$ measures the supply of independent polynomial functions
whose sizes in the weighted box are at most a constant times $X^T$.

\begin{theorem}\label{thm:affine-counting}
Fix integers $1\le r\le N$, a degree bound $e_0$, a constant $C$, positive coordinate
weights $\omega_l$, and positive numbers $\delta_1,\ldots,\delta_r$.
For $X\ge2$, let $S\subseteq\Z^N$ satisfy
\[
 |z_l|\le C X^{\omega_l}\quad(z\in S),
\]
and suppose $S$ lies on an algebraic set $Z$ of degree at most $e_0$
and dimension at most $r$.  Assume the following uniform alternative.
For every fixed degree bound $e$, each irreducible subvariety
$A\subseteq Z$ defined over $\Q$, of degree at most $e$ and dimension
$h>0$, meeting $S$, satisfies either
\begin{equation}\label{eq:affine-exception}
 |S\cap A|\ll_{e,\alpha} X^\alpha
 \quad\text{for every fixed }\alpha>0,
\end{equation}
or
\begin{equation}\label{eq:affine-hilbert-input}
 H_A(T)\ge \frac{T^h}{h!\delta_h^h}-C_eT^{h-1}
 \quad(T\ge1).
\end{equation}
The constants in these alternatives are uniform in $X,Z,S,A$.
Put $E_h=\max_{h\le s\le r}\delta_s$.  Then, for every $\eps>0$,
\begin{equation}\label{eq:affine-counting-conclusion}
 |S|\ll_\eps X^{E_1+\cdots+E_r+\eps}.
\end{equation}
The implied constant depends only on the fixed data and the uniform
constants in the hypotheses.  In particular, no bound on the
coefficients initially defining $Z$ is required.
\end{theorem}

The coordinate weights are fixed in this statement.  Uniformity over a
family of weights follows from the proof when their positive lower
bounds, upper bounds, and the constants in
\eqref{eq:affine-hilbert-input} are uniform.  Empty sets cause no issue;
zero-dimensional sets have at most their degree many points.  We may
therefore assume $1\le r\le N$ and $S\ne\varnothing$.

\subsection{Smooth patches with controlled equations}

A \emph{patch} will mean a triple $(A,G,\Delta)$, where $A$ is an
irreducible $\Q$-defined variety of dimension $h$, the list $G$ consists
of $N-h$ integer polynomials vanishing on $A$, and $\Delta$ is a square
Jacobian minor in $N-h$ chosen coordinate columns.  It covers the points
of $A$ at which $\Delta\ne0$.  If $h=N$, the list is empty and
$\Delta=1$.  A patch need not describe $A$ globally as a complete
intersection: its equations provide local coordinates wherever its
minor is nonzero.

\begin{lemma}\label{lem:affine-patches}
Let $S$ satisfy the box bound in Theorem~\ref{thm:affine-counting}.
For every algebraic set $Z'\subseteq\C^N$ of degree at most $e$, there
is a list of boundedly many patches $(A,G,\Delta)$ with $A\subseteq Z'$
that covers $S\cap Z'$.  Their degrees are bounded, and all coefficients
of $G$ and $\Delta$ are bounded by $C'X^{C'}$, where the bounds depend
only on $e,N,C$ and the weights.  In particular, at a covered point
$z$, the integer $\Delta(z)$ is nonzero and has absolute value at most
$C''X^{C''}$ with the same uniform dependence.
The construction can be fixed as a function of $Z'$ and $S$ alone.
\end{lemma}

\begin{proof}
Split $Z'$ into irreducible components.  If a component $A$ is not
defined over $\Q$, some automorphism $\sigma$ of $\C$ fixing $\Q$
satisfies $\sigma(A)\ne A$.  Every rational point of $A$ lies in
$A\cap\sigma(A)$, a proper intersection of bounded degree.  We replace
$A$ by the components of this intersection and continue in smaller
dimension.

Here is a justification of the descent assertion used in this step.
If all such automorphisms fixed $A$, they would preserve each
finite-dimensional space of its polynomial equations of degree at
most $e'$.  The reduced row-echelon basis in the ordered monomial
coordinates is unique, so its coefficients would be fixed by all these
automorphisms and hence rational.  These spaces would give rational
equations for $A$.  The fixed field is indeed $\Q$: a nonrational
algebraic number has a conjugate, while a transcendental number can be
included in a transcendence basis over $\overline\Q$ and moved, for
example by translation in that basis; these maps extend to
automorphisms of $\C$.

Now let $A$ be defined over $\Q$, of dimension $h$ and degree $e'$.
The equations of ordinary degree at most $e'$ in its ideal have generic
Jacobian rank $N-h$.  To see this, choose $h$ coordinate functions
algebraically independent in $\C(A)$.  Project onto these coordinates
and, in turn, each remaining coordinate.  Each image closure is an
irreducible hypersurface of degree at most $e'$.  Its equation has
nonzero derivative in the extra coordinate, by separability in
characteristic zero.  Lifting these equations to $\C^N$ gives a
diagonal, generically nonzero minor in the remaining $N-h$ coordinates.

Evaluate all monomials of degree at most $e'$ on the finite set
$S\cap A$.  The resulting integer matrix has a bounded number of
columns, all entries bounded by a fixed power of $X$.  Its nullspace
$W$ has an integer spanning set of polynomial height: select independent
rows and pivot columns, solve for the nonpivot variables by cofactors,
and clear the pivot determinant.  All determinants have bounded size.
We compare its complexification with the degree-$e'$ part of the ideal
of $A$, which it contains.

If this containment is strict, one bounded-height spanning polynomial
of $W$ does not vanish identically on $A$.  Its zero set contains
$S\cap A$, so we replace $A$ by the components of the resulting proper
intersection and recurse.  If equality holds, the bounded-height
spanning equations furnish the generic Jacobian rank just proved.
Choose $N-h$ of them and a minor $\Delta$ nonzero on $A$.  They give
one patch, and we recurse on $A\cap\{\Delta=0\}$ for its uncovered
points.  A rational zero-dimensional component containing $z\in S$
is the point $z$; use the equations $z_l'-z_l$ and minor $1$.

Every recursive intersection is proper, with degree bounded in terms
of the preceding degree.  The dimension decreases along every branch,
so B\'ezout's bounds control the length of the list and all its degrees.
The same argument controls all height exponents.  Finally, fix the
choices and ordering once as functions of the input algebraic set,
its fixed degree allowance, and $S$.  Equal inputs then give identical
lists, equations, minors and orderings.
\end{proof}

The last convention is essential below.  A prime used to define an
input variety may affect that variety, but it is not separately
consulted when choosing its patch list.

\subsection{A determinant cut in a simultaneous residue class}

For $h\ge1$, let $s_h(L)$ be the sum of the first $L$ total degrees
of monomials in $h$ variables, ordered by degree and counted with their
monomial multiplicities.  Counting lattice points in a simplex gives
\begin{equation}\label{eq:affine-local-orders}
 s_h(L)=\frac{h}{h+1}(h!)^{1/h}L^{1+1/h}+O_h(L).
\end{equation}

\begin{lemma}\label{lem:affine-residue-cut}
Fix $\nu>0$ and a bound on the degree of a patch $(A,G,\Delta)$ and
its equations, with $h=\dim A>0$.  Suppose $A$ satisfies
\eqref{eq:affine-hilbert-input}.  Let $z\in S\cap A$ be covered by the
patch, and let $\mathcal Q$ be distinct primes such that
\[
 q\nmid\Delta(z)\quad(q\in\mathcal Q),\qquad
 \prod_{q\in\mathcal Q}q\ge X^{\delta_h+\nu}.
\]
For sufficiently large $X$, the set of all points in $S\cap A$
congruent to $z$ modulo these primes lies on a polynomial of bounded
degree that does not vanish identically on $A$.  The degree bound is
uniform in the patch and primes, with dependence on $\nu$ and the
fixed data.  Additional congruence conditions may also be imposed.
\end{lemma}

\begin{proof}
Fix $q\nmid\Delta(z)$.  Translate coordinates by $z$, taking as free
coordinates the $h$ columns outside the chosen minor.  The implicit
equations $G=0$ solve the other coordinates as power series in these
free differences, with coefficients in $\Z_q$: recursively solve by
the Jacobian matrix, which is invertible over $\Z_q$.  The series
converge on $q\Z_q^h$, and Hensel uniqueness gives every solution to
$G=0$ congruent to $z$.  Hence each coordinate monomial has such an
integral expansion on the entire residue class.

An $L$-by-$L$ determinant of monomial values at points in that class
has $q$-adic valuation at least $s_h(L)$.  Indeed, expand in monomials
of the free differences.  A term with a repeated expansion monomial
vanishes, and the differences belong to $q\Z_q$.  The least total
degree of $L$ distinct expansion monomials is $s_h(L)$.  Truncation
modulo arbitrary powers of $q$ justifies this argument for the
convergent series.

On the other hand, \eqref{eq:affine-hilbert-input} permits us to choose
independent coordinate monomials successively, with the $j$-th of
weight at most
\[
 \delta_h(h!j)^{1/h}+O(1).
\]
To obtain this bound, substitute
$T=\delta_h(h!j)^{1/h}+C_1$ in the Hilbert lower bound and choose the
fixed $C_1$ sufficiently large.  Summing the weights and using
\eqref{eq:affine-local-orders}, the first $L$ chosen monomials have
total weight at most $\delta_h s_h(L)+O(L)$.

The absolute value of their evaluation determinant is therefore at
most
\[
 C_L X^{\delta_h s_h(L)+C_2L}.
\]
It is an integer divisible by
$\bigl(\prod_{q\in\mathcal Q}q\bigr)^{s_h(L)}$.
Choose $L$ sufficiently large that $\nu s_h(L)>C_2L$, and then take
$X$ large.  Every such determinant is zero.  Thus the evaluations of
these independent monomials on the full class have deficient rank.
A nonzero rational linear combination vanishes on the class but not
on $A$.  Its ordinary degree is bounded because its weight and the
positive minimum coordinate weight are bounded.  All these choices
are uniform under the stated hypotheses.
\end{proof}

We have obtained polynomial cuts whenever a sufficiently large product
of primes is good for the current patch.  The remaining issue is that
the patch itself will depend on earlier primes.  We address this by
recording the entire finite sequence of evaluation ranks.

\subsection{Adaptive paths and a deletion estimate}

To compare a block of primes with the same block after one prime is
deleted, we must define the successive patches for both choices before
knowing which primes are good. We therefore construct all $r$ stages
for every point and every block choice, even when the expected dimension
drops fail. At each stage we retain the whole space of bounded-degree
polynomials vanishing on the current residue class. Equal-dimensional
nested spaces are equal; this will make unchanged evaluation ranks
force the same subsequent patches. Choosing just one vanishing
polynomial would not provide that implication.

Fix a small $\gamma>0$ and set $h_i=r+1-i$ for $1\le i\le r$.
Choose successive ordinary degree limits $D_i$ for the polynomial
cuts.  At step $i$, take $D_i$ large enough for
Lemma~\ref{lem:affine-residue-cut} in parent dimension $h_i$, with
$\nu=\gamma/2$.  This can be done inductively.  Intersecting with the
common zero set of a polynomial space of degree at most $D_i$, then
applying Lemma~\ref{lem:affine-patches}, gives degree bounds determined
by the parent bound and $D_i$.  These bounds hold even when dimension
does not decrease.  Consequently they cover every path of the fixed
length $r$.

Let $\mathcal B_1,\ldots,\mathcal B_r$ be disjoint finite blocks of
primes, and write
\[
 Q_i=\prod_{j\le i}\prod_{q\in\mathcal B_j}q.
\]
For each $z\in S$, define a path through all $r$ steps as follows.
Take the first patch $(A_0,G_0,\Delta_0)$ of $Z$ covering $z$.
Given its parent patch at step $i$, let
\begin{equation}\label{eq:affine-evaluation-kernel}
 \mathcal K_i=
 \left\{F\in\Q[z_1,\ldots,z_N]_{\le D_i}:
 F(z')=0\text{ for all }z'\in S\cap A_{i-1},\ z'\equiv z\pmod{Q_i}
 \right\}.
\end{equation}
Record its codimension $\rho_i$ in the polynomial space.  Take the
first patch $(A_i,G_i,\Delta_i)$ covering $z$ in
$A_{i-1}\cap V(\mathcal K_i)$.  Every such patch exists because the
intersection contains $z$.  In particular, all $\Delta_i(z)$ are
nonzero integers of uniformly polynomial size. An exceptional
alternative or a zero-dimensional parent does not stop this
construction.

Adding a prime to one block makes the vector
$(\rho_1,\ldots,\rho_r)$ lexicographically nonincreasing.  Until its
first changed entry the parent patches agree: the new residue class
is a subset of the old one, so the kernels are nested, and equal
codimensions make them equal.  Their intersection sets and deterministic
patch lists therefore agree as well.  The first changed codimension
can only decrease.  Moreover, equality of the complete rank vectors
forces equality of all kernels and patches along the paths.

Put $M_i=\binom{N+D_i}{N}$.  Encode the rank vector by the integer
\begin{equation}\label{eq:affine-rank-record}
 \mathcal R=\sum_{i=1}^r\rho_i\prod_{j>i}(M_j+1),
 \qquad 0\le\mathcal R<B_0:=\prod_{i=1}^r(M_i+1).
\end{equation}
It is nonincreasing when a prime is added.  An unchanged encoding
gives identical evaluated minors along the entire path.

\begin{lemma}[Deletion estimate]\label{lem:affine-deletion}
Fix $z\in S$, all blocks except the $i$-th, and a prime pool
$\mathcal P_i$ disjoint from those blocks, whose primes lie between
$X^\zeta$ and $C_rX^\zeta$,
where $\zeta>0$ is fixed and $|\mathcal P_i|\longrightarrow\infty$.
Choose $t$ uniformly among $T_0$ consecutive positive integers
$\ell,\ldots,\ell+T_0-1$, fixed independently of $X$, and then choose
$\mathcal B_i$ uniformly among the $t$-subsets of the pool.  The
expected fraction of primes in $\mathcal B_i$ dividing at least one
of the full-path integers $\Delta_0(z),\ldots,\Delta_r(z)$ is at most
\[
 B_0/T_0+o(1).
\]
The error is uniform in $z$ and in the fixed other blocks.
\end{lemma}

\begin{proof}
Delete a uniformly chosen prime $q$ from $\mathcal B_i$.  Let $F(t)$
be the expected record \eqref{eq:affine-rank-record} for a uniform
$t$-subset.  The deleted subset is uniform of size $t-1$, so the
expected increase in the record is $F(t-1)-F(t)$.  This increase is
a nonnegative integer and is at least one whenever the record
changes.  Averaging over the consecutive sizes gives
\begin{equation}\label{eq:affine-deletion-change}
 \mathbb P(\text{record changes})
 \le\frac{F(\ell-1)-F(\ell+T_0-1)}{T_0}
 \le\frac{B_0}{T_0}.
\end{equation}

Condition now on the deleted subset and all other blocks.  The
deleted-path minors at $z$ are fixed nonzero integers of uniformly
polynomial size.  Their product has only $O(1/\zeta)$ prime divisors
at least $X^\zeta$.  The removed prime is uniform in the complement
of the deleted subset in $\mathcal P_i$, whose size tends to infinity.
Its conditional probability of dividing that product is therefore
$o(1)$, uniformly.

No conditioning on record stability is made in this last assertion.
Rather, if the record is unchanged, the full and deleted paths have
identical minors.  Consequently
\[
 \{q\text{ divides a full-path minor}\}
 \subseteq
 \{\text{record changes}\}
 \cup\{q\text{ divides a deleted-path minor}\}.
\]
The probability is bounded by \eqref{eq:affine-deletion-change} plus
$o(1)$.  Choosing $q$ uniformly from the full block samples exactly
its fraction of primes dividing a full-path minor, proving the lemma.
\end{proof}

\subsection{Good cumulative products and the counting tree}

We now choose blocks for which at least half of $S$ has enough good
primes at every stage. As the scheduled dimension falls from $r$ to
$1$, the determinant thresholds are $\delta_r,\ldots,\delta_1$.
Earlier primes remain available at later stages, so we allocate the
nondecreasing cumulative exponents
$E_r,\ldots,E_1$ instead. Recall that
$E_h=\max_{h\le s\le r}\delta_s$. Put $E_{r+1}=0$ and
\[
 a_i=E_{h_i}-E_{h_i+1}\ge0.
\]
Choose a small $\xi>0$ such that
\begin{equation}\label{eq:affine-bad-fraction}
 \xi(E_1+r\gamma)<\gamma/2.
\end{equation}
Take $T_0$ large in terms of $r,B_0,\xi$, then take $\zeta>0$ so small
that $\zeta(T_0+1)\le\gamma$.  Use disjoint prime pools of cardinality
tending to infinity, all between $X^\zeta$ and $C_rX^\zeta$.
The lower and upper Chebyshev bounds for the prime-counting function
give such pools for sufficiently large fixed $C_r$: partition the
primes in that interval into $r$ pools of comparable size.
Choose the blocks independently, with block $i$ sampled as in
Lemma~\ref{lem:affine-deletion}, starting its size interval at
\[
 \ell_i=\left\lceil\frac{a_i+\gamma}{\zeta}\right\rceil.
\]
Every allowed size $t_i$ then satisfies
\begin{equation}\label{eq:affine-block-sizes}
 a_i+\gamma\le\zeta t_i\le a_i+2\gamma.
\end{equation}
All sizes are constants as $X$ tends to infinity.

Call $z\in S$ \emph{regular} if in each block at most a fraction $\xi$
of its primes divide any of the full-path minors at $z$.
Lemma~\ref{lem:affine-deletion}, Markov's inequality and a union bound
give
\[
 \mathbb P(z\text{ is not regular})
 \le \frac r\xi\bigl(B_0/T_0+o(1)\bigr)<\frac12
\]
for our sufficiently large $T_0$ and then sufficiently large $X$.
Summing over the finite set $S$, some fixed choice of blocks makes at
least $|S|/2$ points regular.  Fix those blocks.

For a regular point, all but a fraction $\xi$ of the primes in each
of the first $i$ blocks avoid the current parent minor
$\Delta_{i-1}(z)$.  Their product is at least
\begin{align}
 X^{(1-\xi)\zeta\sum_{j\le i}t_j}
 &\ge X^{(1-\xi)(E_{h_i}+i\gamma)}
 \ge X^{\delta_{h_i}+\gamma/2}.
 \label{eq:affine-cumulative-product}
\end{align}
The first inequality uses
$\sum_{j\le i}a_j=E_{h_i}$; the last uses
\eqref{eq:affine-bad-fraction}.  This cumulative product includes
primes chosen before the current parent was constructed.

The full paths were needed to compare prime choices in the deletion
estimate. We now arrange the full paths of regular points in a tree,
using the same kernels and fixed patch lists, and stop whenever a bound
is already available. Start from their bounded list of initial patches;
a node records its parent patch and the previous residue classes.
Stop at a parent satisfying
\eqref{eq:affine-exception}, charging all its points by that bound;
also stop at a zero-dimensional parent.  Otherwise continue as
follows.  We show inductively that a parent at step $i$ has dimension
at most $h_i$.

Fix such a parent patch $(A,G,\Delta)$ of dimension $h\le h_i$.
At an auxiliary prime $q$, its reductions with $\Delta\ne0$ number
$O(q^h)$.  Indeed, fix the $h$ free coordinates determined by the
minor.  The other $N-h$ equations have invertible Jacobian at every
point being counted, so these points are isolated zeros of that
square system.  B\'ezout bounds their number by a fixed constant,
even if other components of the fiber have positive dimension.
Summing over the free coordinates proves the bound.  It uses no
assumption that the whole variety has good reduction.

For regular points, the primes of the new block at which this parent
minor vanishes form a subset of size at most $\xi t_i$.  Allow all
such subsets; their number is at most $2^{t_i}$, a constant.  At those
primes use the trivial $q^N$ bound and at the others use $O(q^{h_i})$.
Since $q\le C_rX^\zeta$, the number of new simultaneous residue
classes is at most
\begin{equation}\label{eq:affine-branch-count}
 O\!\left(X^{\zeta t_i(h_i+N\xi)}\right)
 \le O\!\left(X^{(a_i+2\gamma)(h_i+N\xi)}\right).
\end{equation}
The constants from B\'ezout raised to $t_i$ are still independent of
$X$.  Imposing earlier congruences can only decrease this count.

Consider one resulting class containing a regular point $z$.  If the
parent has dimension exactly $h_i$, then
\eqref{eq:affine-cumulative-product} and
Lemma~\ref{lem:affine-residue-cut} give a polynomial in $\mathcal K_i$
that does not vanish identically on the parent.  This polynomial
vanishes on the \emph{entire} class in $S\cap A$: every class member
has the same residues as $z$ at the good primes and satisfies $G=0$.
Thus it is a cut for the kernel in
\eqref{eq:affine-evaluation-kernel}, not merely for its regular points.
Consequently $A\cap V(\mathcal K_i)$ has dimension at most $h_i-1$,
and its fixed patch list has boundedly many members. If the parent had
smaller dimension already, its children automatically satisfy the same
bound. This proves the dimension induction.

After $r$ stages every nonexceptional leaf is zero-dimensional and
contains boundedly many points.  The total branching exponent from
\eqref{eq:affine-branch-count} is at most
\begin{equation}\label{eq:affine-exponent-accounting}
 \sum_{i=1}^r h_i a_i
 +2\gamma\sum_{i=1}^r h_i
 +N\xi(E_1+2r\gamma).
\end{equation}
The first term telescopes to
\[
 \sum_{h=1}^r h(E_h-E_{h+1})=\sum_{h=1}^r E_h.
\]
Every prefix of the tree has no larger exponent, because all the
block increments and errors are nonnegative.  Hence stopping at
exceptional parents multiplies this bound only by $O(X^\alpha)$,
using \eqref{eq:affine-exception}; the bounded number of possible
stopping depths contributes a constant.

To complete the choices for a prescribed $\eps$, first take $\gamma$
so small that the second term of
\eqref{eq:affine-exponent-accounting} is less than $\eps/4$.
Choose the degree limits and hence $B_0$.  Then take $\xi$ small
enough for \eqref{eq:affine-bad-fraction} and for the last term of
\eqref{eq:affine-exponent-accounting} to be less than $\eps/4$.
Take $\alpha<\eps/4$, then choose $T_0$, $\zeta$ and finally $X$
as above.  This order makes every degree, height and probability
constant fixed before $X$ tends to infinity.  We have bounded at
least half of $S$ by the right side of
\eqref{eq:affine-counting-conclusion}; multiplying by two proves
Theorem~\ref{thm:affine-counting}.
